This section records ensemble extensions that are immediate from the principle,
and separates them from extensions that require new input.

\subsection{Division algebras}

The real and complex Gaussian cases correspond to real dimensions \(d=1\) and
\(d=2\) of one scalar coordinate.  The same hard-edge bookkeeping applies to
quaternionic Gaussian matrices, where one scalar has real dimension \(d=4\).

For quaternionic matrices we use the following conventions; see
\cite{Zhang1997Quaternions} for quaternion matrix background.  Row vectors in
\(\HH^m\) carry the inner product
\[
  \langle x,y\rangle_{\HH}=\sum_{j=1}^m x_j\overline{y_j},
\]
linear in the first variable, and \(\spn_{\HH}\) denotes left
\(\HH\)-span.  For \(A\in\HH^{n\times m}\),
\[
  \smin(A)=\inf_{\|x\|_2=1}\|Ax\|_2 .
\]
Equivalently, \(\smin(A)^2\) is the least eigenvalue of \(A^*A\), and it is
also the least singular value of the realification; the realification only
repeats quaternionic singular values with multiplicity four.  A standard
quaternion Gaussian scalar is
\[
  q=q_0+q_1\mathbf i+q_2\mathbf j+q_3\mathbf k,
  \qquad q_\ell\sim N(0,1/4)\quad\hbox{iid},
\]
so \(\EE|q|^2=1\), \(|q|^2=(1/4)\chi^2_4\), and
\[
  p_{\HH}(t):=\PP\{|q|\le t\}=2t^4+O(t^6)\qquad(t\downarrow0).
\]
The compact symplectic group is
\[
  \Sp(m)=\{U\in\HH^{m\times m}:U^*U=UU^*=I_m\},
\]
acting on row vectors by right multiplication.

\begin{theorem}[Quaternion Gaussian extremal hard edge]\label{thm:quaternion-gaussian-base}
Let \(A_m\in\HH^{N_m\times m}\) have iid standard quaternion Gaussian entries
and \(N_m/m\to\gamma>1\).  Put
\[
  M_m^{\HH}=\min_{\substack{T\subset[N_m]\\ |T|=m}}\smin(A_{m,T}).
\]
Then
\[
  \frac{1}{m}\log M_m^{\HH}
  \xrightarrow{\PP}-\frac{h(\gamma)}{4}.
\]
\end{theorem}

\begin{proof}
The lower side is a quaternionic one-matrix hard-edge bound.  Let
\(G\in\HH^{m\times m}\) have iid standard quaternion Gaussian rows \(a_i\).
The quaternionic negative second-moment identity gives
\[
  \sum_{k=1}^m\sigma_k(G)^{-2}
  =
  \sum_{i=1}^m
  \dist\big(a_i,\spn_{\HH}\{a_j:j\ne i\}\big)^{-2}.
\]
Thus \(\smin(G)\le t\) forces, for some \(i\),
\[
  \dist\big(a_i,\spn_{\HH}\{a_j:j\ne i\}\big)\le \sqrt m\,t .
\]
Conditioning on the other rows, this distance is
\(|\langle a_i,u\rangle_{\HH}|\) for an \(\HH\)-unit normal \(u\), and this
scalar is a standard quaternion Gaussian.  Hence, for \(0<t\le1\),
\[
  \PP\{\smin(G)\le t\}
  \le m\,p_{\HH}(\sqrt m\,t)
  \le C m^3t^4.
\]
The union bound over \(\binom{N_m}{m}=e^{(h(\gamma)+o(1))m}\) row sets gives
\[
  \PP\{M_m^{\HH}\le e^{-sm}\}\to0\qquad(s>h(\gamma)/4).
\]

For the upper side, fix \(s<h(\gamma)/4\), set \(t=e^{-sm}\), and write
\(N=N_m\).  Let \(x=a_N\).  For each
\((m-1)\)-set \(S\subset[N-1]\), let \(L_S\) be the quaternionic normal line to
\(\spn_{\HH}\{a_i:i\in S\}\), and define
\[
  E_S=\{\dist(x,\spn_{\HH}\{a_i:i\in S\})\le t\}.
\]
Equivalently, for any unit representative \(u_S\in L_S\),
\(E_S=\{|\langle x,u_S\rangle_{\HH}|\le t\}\); this condition is independent of
the unit-quaternion phase of \(u_S\).  Let
\[
  Z_t=\sum_{\substack{S\subset[N-1]\\ |S|=m-1}}\mathbf1_{E_S},
  \qquad B_m=\binom{N-1}{m-1}.
\]
If \(Z_t>0\), then \(M_m^{\HH}\le t\).  Since \(x\) is independent of the
normal lines,
\[
  \EE Z_t=B_m p_{\HH}(t)=\exp\{(h(\gamma)-4s+o(1))m\}\to\infty.
\]

It remains to show \(\EE Z_t^2/(\EE Z_t)^2\to1\).  Fix two sets \(S,T\) and put
\[
  D=m-|S\cap T|.
\]
The case \(D=1\) is the diagonal \(S=T\).  If \(D\ge2\), condition on the rows
indexed by \(S\cap T\).  Their orthogonal complement is isomorphic to
\(\HH^D\).  The remaining rows in \(S\setminus T\) and \(T\setminus S\) are
independent \(\Sp(D)\)-invariant quaternion Gaussian vectors, so the normal
lines \(L_S,L_T\) are independent Haar points of quaternionic projective space
\(\HH P^{D-1}\).

Choose unit representatives \(u,v\) and set
\(\rho=\langle u,v\rangle_{\HH}\).  The quantities below depend only on
\(|\rho|\).  If \(D\ge3\), then
\[
  |\rho|^2\sim\operatorname{Beta}(2,2(D-1))
\]
and
\[
  \EE\big[(1-|\rho|^2)^{-2}\mid D\big]
  =
  \frac{(2D-1)(2D-2)}{(2D-3)(2D-4)}
  =1+\frac2D+O(D^{-2}).
\]
This negative moment is finite exactly in the range \(D\ge3\), which is why
the case \(D=2\) is split off below.

For fixed \(u,v\) with \(|\rho|<1\), a quaternionic unitary rotation and a
unit-quaternion phase change put
\[
  u=e_1,\qquad
  v=e_1\rho+e_2(1-|\rho|^2)^{1/2}.
\]
In the two quaternionic coordinates \((x_1,x_2)\), the real Jacobian of
\[
  (x_1,x_2)\mapsto
  \big(\langle x,u\rangle_{\HH},\langle x,v\rangle_{\HH}\big)
\]
is \((1-|\rho|^2)^2\).  Since the joint quaternion Gaussian density is maximal
at the origin,
\[
  \PP\{|\langle x,u\rangle_{\HH}|\le t,\,
        |\langle x,v\rangle_{\HH}|\le t\mid u,v\}
  \le (1+O(t^2))p_{\HH}(t)^2(1-|\rho|^2)^{-2}.
\]
For \(D=2\) we only use the trivial bound by \(p_{\HH}(t)\).

For a uniformly chosen ordered pair \((S,T)\) of \((m-1)\)-sets in \([N-1]\),
\[
  \PP\{D=d\}
  =
  \frac{\binom{m-1}{d-1}\binom{N-m}{d-1}}{\binom{N-1}{m-1}},
  \qquad 1\le d\le m,
\]
and
\[
  \EE\frac1D
  =
  \frac{N}{m(N-m+1)}
  =O(m^{-1}).
\]
Also
\[
  \PP\{D=2\}
  =
  \frac{(m-1)(N-m)}{\binom{N-1}{m-1}}.
\]
Therefore
\[
  \frac{\PP\{D=2\}}{p_{\HH}(t)}
  \le m^{O(1)}\exp\{-(h(\gamma)-4s+o(1))m\}\to0,
\]
and the same bound with \(\PP\{D=1\}=1/B_m\) handles the diagonal contribution.
Combining the estimates,
\[
  \frac{\EE Z_t^2}{(\EE Z_t)^2}
  \le o(1)
     +(1+O(t^2))
      \sum_{d=3}^m \PP\{D=d\}\bigl(1+2/d+O(d^{-2})\bigr)
  =1+o(1).
\]
The reverse inequality is automatic from \(\Var Z_t\ge0\).  Chebyshev's
inequality gives \(\PP\{Z_t=0\}\to0\), hence
\(\PP\{M_m^{\HH}\le e^{-sm}\}\to1\) for every \(s<h(\gamma)/4\).  Together
with the lower side, this proves the claimed convergence.
\end{proof}

\begin{theorem}[Quaternion random number of Gaussian rows]\label{thm:quat-random-row-count-gaussian}
For each \(m\), let \((R_{im})_{i\ge1}\) be an infinite sequence of independent
standard quaternion Gaussian rows in \(\HH^m\).  For an integer \(n\ge m\), put
\[
  M_m^{\HH}(n)=
  \min_{\substack{T\subset[n]\\ |T|=m}}\smin((R_{im})_{i\in T}).
\]
Let \(N_m\) be integer-valued random variables, possibly dependent on the row
array, such that
\[
  \frac{N_m}{m}\xrightarrow{\PP}\gamma>1 .
\]
Extend \(M_m^{\HH}(N_m)\) by any fixed positive value on the event \(N_m<m\).
This event has probability tending to zero, and
\[
  \frac1m\log M_m^{\HH}(N_m)
  \xrightarrow{\PP}
  -\frac{h(\gamma)}4 .
\]
The normalized square Gram statistic
\[
  \Gamma_m^{\HH}(N_m)=
  \min_{\substack{T\subset[N_m]\\ |T|=m}}
  \lambda_{\min}\!\left(m^{-1}R_T^*R_T\right)
\]
has logarithmic limit \(-h(\gamma)/2\), with the same convention on
\(\{N_m<m\}\).  If the row arrays are realized on a common probability space
and \(N_m/m\to\gamma\) almost surely, then the corresponding lower envelopes
also hold:
\[
  \liminf_{m\to\infty}\frac1m\log M_m^{\HH}(N_m)
  \ge -\frac{h(\gamma)}4,\qquad
  \liminf_{m\to\infty}\frac1m\log \Gamma_m^{\HH}(N_m)
  \ge -\frac{h(\gamma)}2
  \qquad\hbox{a.s.}
\]
\end{theorem}

\begin{proof}
The map \(n\mapsto M_m^{\HH}(n)\) is nonincreasing, since increasing the number
of available prefix rows only enlarges the family of tested square submatrices.
Fix
\[
  1<\gamma_-<\gamma<\gamma_+ .
\]
With probability tending to one,
\[
  \lfloor\gamma_-m\rfloor\le N_m\le\lceil\gamma_+m\rceil .
\]
On this event,
\[
  M_m^{\HH}(\lceil\gamma_+m\rceil)
  \le M_m^{\HH}(N_m)
  \le M_m^{\HH}(\lfloor\gamma_-m\rfloor).
\]
Theorem~\ref{thm:quaternion-gaussian-base}, applied to the two deterministic
row counts, gives limits \(-h(\gamma_+)/4\) and \(-h(\gamma_-)/4\).  Letting
\(\gamma_-\uparrow\gamma\) and \(\gamma_+\downarrow\gamma\) proves the
convergence in probability.  No independence between \(N_m\) and the row array
is used for this prefix-count statement.

The Gram convergence follows from Corollary~\ref{cor:gram-transfer}.  For the
almost-sure lower envelope, fix \(\gamma_+>\gamma\).  Eventually almost surely,
\(N_m\le\lceil\gamma_+m\rceil\), hence
\[
  M_m^{\HH}(N_m)\ge M_m^{\HH}(\lceil\gamma_+m\rceil).
\]
The summable lower-side estimate in the proof of
Theorem~\ref{thm:quaternion-gaussian-base} gives the deterministic-row-count
lower envelope at \(\gamma_+\); then let \(\gamma_+\downarrow\gamma\).  The
Gram lower envelope follows by the same pathwise squaring identity.
\end{proof}

\begin{corollary}[Quaternion Poissonization and independent thinning]\label{cor:quat-poisson-thinning}
The conclusion of Theorem~\ref{thm:quat-random-row-count-gaussian} applies when
\(N_m\) is independent of the quaternion Gaussian row values and one of the
following conditions holds:
\begin{enumerate}
\item \(N_m\) is Poisson with mean \(\nu_m\), and \(\nu_m/m\to\gamma>1\);
\item \(N_m\) is binomial \(\operatorname{Bin}(n_m^{\rm pool},p_m)\), with
      \(n_m^{\rm pool}p_m/m\to\gamma>1\);
\item \(N_m=\sum_{i=1}^{n_m^{\rm pool}}\xi_{im}\), where the \(\xi_{im}\) are
      independent Bernoulli retention indicators, independent of the row
      values, and \(\nu_m=\sum_i\EE\xi_{im}\) satisfies
      \(\nu_m/m\to\gamma>1\).
\end{enumerate}
In the thinning case, conditional on the retained count, the retained rows can
be reindexed as iid quaternion Gaussian rows; the statistic is invariant under
row order.  Value-dependent thinning or selection is not covered.
\end{corollary}

\begin{corollary}[Quaternion random-prefix conditioning transfer]\label{cor:quat-random-prefix-conditioning-exact}
For each \(m\), let \((R_{im})_{i\ge1}\) be an infinite sequence of independent
standard quaternion Gaussian rows in \(\HH^m\).  Let \(N_m\) be integer-valued,
possibly dependent on the row array, with
\[
  \frac{N_m}{m}\xrightarrow{\PP}\gamma>1 .
\]
Let \(d_{im}>0\), \(i\ge1\), be finite real random row weights and let
\(S_m\in\HH^{m\times m}\) be invertible.  The weights and \(S_m\) may be
random and dependent on \(N_m\) and on the row array.  For
\(T\subset[N_m]\), \(|T|=m\), set
\[
  B_{m,T}=\operatorname{diag}(d_{im}:i\in T)\,(R_{im})_{i\in T}\,S_m .
\]
Assume
\[
  \Delta_m^{\rm rpref}:=
  \max_{1\le i\le N_m}|\log d_{im}|
  +\bigl|\log\smin(S_m)\bigr|+\bigl|\log\|S_m\|\bigr|
  =o_\PP(m).
\]
Absent row weights or the absent right factor are omitted.  On the event
\(N_m<m\), give the statistics below any fixed positive value.  Then
\[
  \frac1m\log
  \min_{\substack{T\subset[N_m]\\ |T|=m}}\smin(B_{m,T})
  \xrightarrow{\PP}
  -\frac{h(\gamma)}4
\]
and
\[
  \frac1m\log
  \min_{\substack{T\subset[N_m]\\ |T|=m}}
  \lambda_{\min}\!\left(m^{-1}B_{m,T}^*B_{m,T}\right)
  \xrightarrow{\PP}
  -\frac{h(\gamma)}2 .
\]
If the row arrays are realized on a common probability space,
\(N_m/m\to\gamma\) almost surely, and
\(\Delta_m^{\rm rpref}=o(m)\) almost surely, then the corresponding lower
envelopes hold:
\[
  \liminf_{m\to\infty}\frac1m\log
  \min_{\substack{T\subset[N_m]\\ |T|=m}}\smin(B_{m,T})
  \ge -\frac{h(\gamma)}4,
\]
and the normalized square Gram lower envelope has base \(-h(\gamma)/2\).

Positive real column diagonal factors are included by taking \(S_m\) diagonal.
Consequently, row-normalized random-prefix quaternion Gaussian matrices and
full-prefix column-normalized quaternion Gaussian matrices have the same
square-row bases, provided the prefix count hypotheses above hold.  For
independent Bernoulli thinning, this is read after the retained iid rows are
reindexed and after the same conditioning-extrema hypothesis is imposed.
\end{corollary}

\begin{proof}
On the event \(N_m\ge m\), for every \(T\subset[N_m]\), \(|T|=m\),
\[
  e^{-\Delta_m^{\rm rpref}}\smin((R_{im})_{i\in T})
  \le
  \smin(B_{m,T})
  \le
  e^{\Delta_m^{\rm rpref}}\smin((R_{im})_{i\in T}).
\]
The inequalities are the usual singular-value bounds for a positive real left
diagonal factor and a global right factor.  They are pathwise and uniform in
\(T\), so after taking minima,
\[
  \left|
  \frac1m\log
  \min_{\substack{T\subset[N_m]\\ |T|=m}}\smin(B_{m,T})
  -
  \frac1m\log M_m^{\HH}(N_m)
  \right|
  \le \frac{\Delta_m^{\rm rpref}}m
\]
outside the event \(N_m<m\), whose probability tends to zero.  The singular
convergence follows from
Theorem~\ref{thm:quat-random-row-count-gaussian}.  The Gram convergence follows
from the pathwise identity in Corollary~\ref{cor:gram-transfer}.  The
almost-sure lower envelopes use the same lower sandwich together with the
almost-sure lower envelopes in Theorem~\ref{thm:quat-random-row-count-gaussian}.

For row normalization, on the high-probability bracket
\(N_m\le\lceil\gamma_+m\rceil\), the row-norm extrema over the realized prefix
are bounded by the extrema over the deterministic prefix
\(\lceil\gamma_+m\rceil\).  For column normalization, on
\(\lfloor\gamma_-m\rfloor\le N_m\le\lceil\gamma_+m\rceil\), each column norm is
bracketed between the corresponding deterministic-prefix column norms.  The
same \(\chi^2\) upper-tail and ball-volume lower-tail estimates used in
Corollary~\ref{cor:quat-square-diagonal-normalized-exact} give
subexponential extrema in probability and, on common couplings with eventual
deterministic brackets, almost surely.
\end{proof}

\begin{corollary}[All-field square random-prefix row and right-factor exact transfer]\label{cor:random-prefix-square-row-right-factor-exact}\label{cor:random-prefix-square-diagonal-exact}
Let \(\F\in\{\R,\C,\HH\}\), with \(d_\F=1,2,4\), respectively.  For each
\(m\), let \((R_{im})_{i\ge1}\) be an infinite sequence of iid standard
\(\F^m\)-Gaussian rows.  Let \(N_m\) be integer-valued, possibly dependent on
the row array, with
\[
  \frac{N_m}{m}\xrightarrow{\PP}\gamma>1 .
\]
Let \(d_{im}>0\), \(i\ge1\), be finite real random row weights and let
\(S_m\in GL_m(\F)\).  The weights and \(S_m\) may be random and dependent on
the rows and on \(N_m\).  Assume
\[
  \Delta_m^{\rm sqrf}
  =
  \max_{1\le i\le N_m}|\log d_{im}|
  +|\log\smin(S_m)|+|\log\|S_m\||
  =o_\PP(m),
\]
with an absent row factor replaced by \(d_{im}\equiv1\) and an absent right
factor replaced by \(S_m=I_m\).  For \(N_m\ge m\) and \(T\subset[N_m]\),
\(|T|=m\), set
\[
  B_{m,T}=\operatorname{diag}(d_{im}:i\in T)\,(R_{im})_{i\in T}\,S_m .
\]
On the event \(N_m<m\), give the following two statistics any fixed positive
finite value.  Then
\[
  \frac1m\log
  \min_{\substack{T\subset[N_m]\\ |T|=m}}\smin(B_{m,T})
  \xrightarrow{\PP}
  -\frac{h(\gamma)}{d_\F},
\]
and
\[
  \frac1m\log
  \min_{\substack{T\subset[N_m]\\ |T|=m}}
  \lambda_{\min}\!\left(m^{-1}B_{m,T}^*B_{m,T}\right)
  \xrightarrow{\PP}
  -\frac{2h(\gamma)}{d_\F}.
\]
If the row arrays are realized on a common probability space,
\(N_m/m\to\gamma\) almost surely, and
\(\Delta_m^{\rm sqrf}=o(m)\) almost surely, then the corresponding lower
envelopes hold:
\[
  \liminf_{m\to\infty}\frac1m\log
  \min_{\substack{T\subset[N_m]\\ |T|=m}}\smin(B_{m,T})
  \ge -\frac{h(\gamma)}{d_\F},
\]
and the normalized square column-Gram lower envelope has base
\(-2h(\gamma)/d_\F\).

Positive real column diagonal factors are included by taking \(S_m\) diagonal.
Consequently, the same conclusions apply to row-normalized random-prefix
Gaussian rows and to full-prefix column-normalized Gaussian columns, where the
column normalization is computed from all rows in the realized prefix \([N_m]\).
\end{corollary}

\begin{proof}
On the event \(N_m\ge m\), for every \(T\subset[N_m]\), \(|T|=m\),
\[
  e^{-\Delta_m^{\rm sqrf}}\smin((R_{im})_{i\in T})
  \le \smin(B_{m,T})
  \le e^{\Delta_m^{\rm sqrf}}\smin((R_{im})_{i\in T}).
\]
The selected rows commute with the single global all-column right factor:
\((R S_m)_T=R_T S_m\).  The displayed estimate is the min--max singular-value
inequality for the positive real row diagonal factor and \(S_m\).  Taking minima
over \(T\) gives
\[
  \left|
  \frac1m\log
  \min_{\substack{T\subset[N_m]\\ |T|=m}}\smin(B_{m,T})
  -
  \frac1m\log M_{\F,m}(N_m)
  \right|
  \le \frac{\Delta_m^{\rm sqrf}}m
\]
outside the event \(N_m<m\), where \(M_{\F,m}(N_m)\) denotes the corresponding
unweighted square-prefix statistic.  The event \(N_m<m\) has probability
tending to zero.  For \(\F=\R,\C\), use
Theorem~\ref{thm:random-row-count-gaussian}; for \(\F=\HH\), use
Theorem~\ref{thm:quat-random-row-count-gaussian}.  The singular-value
convergence follows because \(\Delta_m^{\rm sqrf}=o_\PP(m)\).

The Gram statistic is transferred by the eigenvalue identity
\[
  \lambda_{\min}\!\left(m^{-1}B_{m,T}^*B_{m,T}\right)
  =m^{-1}\smin(B_{m,T})^2.
\]
For the almost-sure lower envelopes, keep only the lower side of the pathwise
sandwich.  In the real and complex cases, the source lower envelope is
Corollary~\ref{cor:random-row-count-bd} with bounded Gaussian density and
\(\kappa=0\); in the quaternion case it is the a.s. part of
Theorem~\ref{thm:quat-random-row-count-gaussian}.  The Gram lower envelope
again follows by squaring.

For row normalization, on the high-probability bracket
\(N_m\le\lceil\gamma_+m\rceil\), the realized-prefix row-norm extrema are bounded
by deterministic-prefix extrema.  For full-prefix column normalization, on
\(\lfloor\gamma_-m\rfloor\le N_m\le\lceil\gamma_+m\rceil\), each column norm is
bracketed between the corresponding deterministic-prefix column norms.  Standard
Gaussian norm upper tails and ball-volume lower tails give
\(\Delta_m^{\rm sqrf}=o_\PP(m)\), and give the a.s. version under common
a.s. prefix-count convergence.

This is only a square random-prefix row/right-factor exact transfer.  It
subsumes Corollary~\ref{cor:quat-random-prefix-conditioning-exact} in the
quaternion case and contains the positive column-diagonal case.  It is the exact square
strengthening of Corollary~\ref{cor:random-prefix-diagonal-fixed-excess-fixed-k}
at \(r=0,k=1\), but it does not extend to fixed-excess exact laws.  Column
selection does not commute with a non-diagonal right factor, so no
selected-column or two-sided non-diagonal covariance statement is implied.  Nor
is there a value-dependent row-selection theorem, selected-submatrix column
normalization, row-Gram statement for \(r>0\), quaternionic non-real row
weighting, condition-number-only covariance claim, or almost-sure exact law.
\end{proof}

\begin{corollary}[Quaternion square diagonal and normalized exact base]\label{cor:quat-square-diagonal-normalized-exact}
Let \(G_m\in\HH^{N_m\times m}\) have iid standard quaternion Gaussian entries
with \(N_m/m\to\gamma>1\).  Let
\[
  \widetilde G_m=D_m^{\rm row}G_mE_m,
\]
where \(D_m^{\rm row}\) and \(E_m\) are positive real diagonal matrices, random
and possibly dependent on \(G_m\), with strictly positive finite entries, and
where either diagonal factor may be absent.  Assume
\[
  \Delta_m^{\rm diag}:=
  \max_i|\log(D_m^{\rm row})_{ii}|
  +
  \max_j|\log(E_m)_{jj}|=o_\PP(m),
\]
omitting absent terms.  Put
\[
  M_{\widetilde G,\HH}
  =
  \min_{\substack{T\subset[N_m]\\ |T|=m}}\smin((\widetilde G_m)_T)
\]
and
\[
  \Gamma_{\widetilde G,\HH}
  =
  \min_{\substack{T\subset[N_m]\\ |T|=m}}
  \lambda_{\min}\!\left(m^{-1}(\widetilde G_m)_T^*
  (\widetilde G_m)_T\right).
\]
Then
\[
  \frac1m\log M_{\widetilde G,\HH}
  \xrightarrow{\PP}
  -\frac{h(\gamma)}4,
  \qquad
  \frac1m\log \Gamma_{\widetilde G,\HH}
  \xrightarrow{\PP}
  -\frac{h(\gamma)}2 .
\]
If \(\Delta_m^{\rm diag}=o(m)\) almost surely on a common coupling, then the
corresponding lower envelopes hold:
\[
  \liminf_{m\to\infty}\frac1m\log M_{\widetilde G,\HH}
  \ge -\frac{h(\gamma)}4,
  \qquad
  \liminf_{m\to\infty}\frac1m\log \Gamma_{\widetilde G,\HH}
  \ge -\frac{h(\gamma)}2
  \qquad\hbox{a.s.}
\]

In particular, the same square-row convergence in probability and lower
envelopes hold for row-normalized quaternion Gaussian rows and for full-matrix
column-normalized quaternion Gaussian columns.
\end{corollary}

\begin{proof}
For every row set \(T\) with \(|T|=m\),
\[
  (\widetilde G_m)_T=(D_m^{\rm row})_T(G_m)_TE_m .
\]
The entries of the diagonal factors are positive real scalars, hence central in
\(\HH\).  The min--max singular-value inequalities therefore give, uniformly
in \(T\),
\[
  e^{-\Delta_m^{\rm diag}}\smin((G_m)_T)
  \le \smin((\widetilde G_m)_T)
  \le e^{\Delta_m^{\rm diag}}\smin((G_m)_T).
\]
Taking the minimum over \(T\), logs, and division by \(m\) gives
\[
  \left|
  \frac1m\log M_{\widetilde G,\HH}
  -\frac1m\log M_{G,\HH}
  \right|
  \le \frac{\Delta_m^{\rm diag}}m
\]
on the almost-sure event that all square row submatrices of \(G_m\) are
nonsingular.  Theorem~\ref{thm:quaternion-gaussian-base} and
\(\Delta_m^{\rm diag}=o_\PP(m)\) give the singular-value convergence.  The
Gram convergence follows from Corollary~\ref{cor:gram-transfer}.

For the almost-sure lower envelope, use the same pathwise comparison together
with the \(r=0,k=1\) lower-envelope case of
Theorem~\ref{thm:quat-fixed-excess-fixed-k-lower} and
Corollary~\ref{cor:quat-rect-kth-gram-lower}.  For row normalization, one row
norm squared is \((1/4)\chi^2_{4m}\); for column normalization, one column norm
squared is \((1/4)\chi^2_{4N_m}\).  The Markov upper-tail and ball-volume
lower-tail estimates used below for quaternion normalization show that the
corresponding logarithmic extrema are \(o_\PP(m)\), and are \(o(m)\) almost
surely after the \(O(m)\) row or column unions.  The rows and columns are
nonzero almost surely.
\end{proof}

\begin{corollary}[Quaternion square row and right-factor conditioning]\label{cor:quat-square-row-covariance-conditioning-exact}
Let \(G_m\in\HH^{N_m\times m}\) have iid standard quaternion Gaussian entries
with \(N_m/m\to\gamma>1\).  Let
\[
  B_m=D_m^{\rm row}G_mS_m,
\]
where \(D_m^{\rm row}\) is a positive real diagonal row factor with strictly
positive finite entries, possibly absent, and \(S_m\in\HH^{m\times m}\) is
invertible.  The factors may be random and dependent on \(G_m\).  Assume
\[
  \Delta_m^{\rm rf}:=
  \max_i|\log(D_m^{\rm row})_{ii}|
  +\bigl|\log\smin(S_m)\bigr|+\bigl|\log\|S_m\|\bigr|=o_\PP(m),
\]
omitting the row term if \(D_m^{\rm row}\) is absent.  Put
\[
  M_{B,\HH}=\min_{\substack{T\subset[N_m]\\ |T|=m}}\smin((B_m)_T)
\]
and
\[
  \Gamma_{B,\HH}
  =
  \min_{\substack{T\subset[N_m]\\ |T|=m}}
  \lambda_{\min}\!\left(m^{-1}(B_m)_T^*(B_m)_T\right).
\]
Then
\[
  \frac1m\log M_{B,\HH}
  \xrightarrow{\PP}
  -\frac{h(\gamma)}4,
  \qquad
  \frac1m\log \Gamma_{B,\HH}
  \xrightarrow{\PP}
  -\frac{h(\gamma)}2 .
\]
If \(\Delta_m^{\rm rf}=o(m)\) almost surely on a common coupling, then the
corresponding lower envelopes hold with the same bases.

In particular, deterministic uniformly elliptic quaternion right factors,
including \(S_m=\Sigma_m^{1/2}\) with
\(0<c\le\lambda_{\min}(\Sigma_m)\le\lambda_{\max}(\Sigma_m)\le C<\infty\),
preserve the square-row quaternion bases.  Absolute singular-value scale is
needed: a subexponential condition number alone would not rule out an
exponential scalar shift.
\end{corollary}

\begin{proof}
For every row set \(T\) with \(|T|=m\),
\[
  (B_m)_T=(D_m^{\rm row})_T(G_m)_TS_m .
\]
Using the quaternionic operator norm and singular-value conventions fixed
above,
\[
  \smin((D_m^{\rm row})_T)\smin(S_m)\smin((G_m)_T)
  \le \smin((B_m)_T)
  \le \|(D_m^{\rm row})_T\|\,\|S_m\|\,\smin((G_m)_T).
\]
The bounds are uniform over \(T\).  Taking minima, logs, and division by \(m\)
shows
\[
  \left|
  \frac1m\log M_{B,\HH}
  -\frac1m\log M_{G,\HH}
  \right|
  \le \frac{\Delta_m^{\rm rf}}m
\]
on the almost-sure event that all square row submatrices of \(G_m\) are
nonsingular.  The singular-value convergence follows from
Theorem~\ref{thm:quaternion-gaussian-base}; the Gram convergence follows from
Corollary~\ref{cor:gram-transfer}.  The lower-envelope statement uses the same
pathwise comparison and the \(r=0,k=1\) lower-envelope cases of
Theorem~\ref{thm:quat-fixed-excess-fixed-k-lower} and
Corollary~\ref{cor:quat-rect-kth-gram-lower}.
\end{proof}

\begin{lemma}[Quaternion fixed-excess fixed-\(k\) Laguerre tail]\label{lem:quat-rect-kth-laguerre-tail}
Fix \(r\ge0\) and \(k\ge1\).  There are constants
\(C_{\HH,r,k}\ge1\) and \(B_{\HH,r,k}\ge0\) such that, for every \(m\ge k\),
every iid standard quaternion Gaussian matrix \(G\in\HH^{(m+r)\times m}\), and
every \(0<t\le1\),
\[
  \PP\{\sigma_{m-k+1}(G)\le t\}
  \le C_{\HH,r,k}m^{B_{\HH,r,k}}t^{4k(k+r)} .
\]
\end{lemma}

\begin{proof}
Let \(0\le\lambda_1\le\cdots\le\lambda_m\) be the eigenvalues of the quaternion
Hermitian matrix \(G^*G\), each counted once as a quaternionic eigenvalue.  With
the standard quaternion Gaussian normalization fixed above, these eigenvalues
have the \(\beta=4\) Laguerre density
\[
  \frac{m!}{Z_m(q)}
  e^{-2\sum_i\lambda_i}
  \prod_{i=1}^m \lambda_i^{q-1}
  \prod_{i<j}(\lambda_j-\lambda_i)^4,
  \qquad q=2(r+1),
\]
in ordered coordinates; this is the quaternion Wishart/Laguerre ensemble, see
\cite[Chs.~3--4]{Forrester2010} together with the quaternion matrix conventions
in \cite{Zhang1997Quaternions}.  Put \(\delta=t^2\).  On
\(\lambda_k\le\delta\), write \(x_i=\lambda_i\) for \(1\le i\le k\) and
\(y_j=\lambda_{k+j}\) for \(1\le j\le m-k\).  Since \(0\le x_i\le y_j\),
\[
  \prod_{i\le k<j}(\lambda_j-\lambda_i)^4
  \le \prod_{j=1}^{m-k}y_j^{4k}.
\]
The integral over the \(k\) bottom variables contributes
\[
  \int_{[0,\delta]^k}\prod_i x_i^{q-1}
       \prod_{i<j}|x_i-x_j|^4\,dx
  =
  C_{r,k}\delta^{kq+4\binom{k}{2}}
  =
  C_{r,k}\delta^{2k(k+r)}.
\]
The remaining variables are bounded by a Laguerre--Selberg integral with
parameter \(q+4k\), so
\[
  \PP\{\lambda_k\le\delta\}
  \le
  C_{r,k}\delta^{2k(k+r)}
  \frac{m!}{(m-k)!}\frac{Z_{m-k}(q+4k)}{Z_m(q)} .
\]
For fixed \(r,k\), the Laguerre--Selberg product and Stirling's formula bound
the ratio by \(C_{r,k}m^{B_{r,k}}\), after absorbing finitely many small \(m\)
into the constant.  Since \(\sigma_{m-k+1}(G)=\sqrt{\lambda_k}\), this gives the
claimed singular-value tail.
\end{proof}

\begin{theorem}[Quaternion fixed-excess fixed-\(k\) lower side]\label{thm:quat-fixed-excess-fixed-k-lower}
Fix \(r\ge0\) and \(k\ge1\).  Let \(A_m\in\HH^{N_m\times m}\) have iid standard
quaternion Gaussian entries, with \(N_m\ge m+r\), and put
\[
  M_{\HH,m}^{(r)}[k]
  =
  \min_{\substack{T\subset[N_m]\\ |T|=m+r}}\sigma_{m-k+1}(A_{m,T}).
\]
There are constants \(C_{\HH,r,k}\ge1\) and \(B_{\HH,r,k}\ge0\) such that, for
every \(u_m>0\),
\[
  \PP\left\{M_{\HH,m}^{(r)}[k]\le
  \exp\left[-\frac{\log\binom{N_m}{m+r}
    +B_{\HH,r,k}\log m+\log C_{\HH,r,k}+u_m}{4k(k+r)}\right]\right\}
  \le e^{-u_m}.
\]
If \(N_m/m\to\gamma>1\), then, for every \(\varepsilon>0\),
\[
  \PP\left\{M_{\HH,m}^{(r)}[k]\le
  \exp\!\left[-\left(\frac{h(\gamma)}{4k(k+r)}+\varepsilon\right)m\right]\right\}
  \to0.
\]
On any common coupling over \(m\),
\[
  \liminf_{m\to\infty}\frac1m\log M_{\HH,m}^{(r)}[k]
  \ge -\frac{h(\gamma)}{4k(k+r)}
  \qquad\hbox{a.s.}
\]
\end{theorem}

\begin{proof}
The finite-budget estimate is the union bound over the
\(\binom{N_m}{m+r}\) selected row sets, using
Lemma~\ref{lem:quat-rect-kth-laguerre-tail}.  Taking
\(u_m=4k(k+r)\varepsilon m\) and absorbing the logarithmic terms gives the
fixed-aspect lower side.  The same choice with, say, half the displayed
\(\varepsilon\)-slack gives summable failure probabilities, so the almost-sure
lower envelope follows from the first Borel--Cantelli lemma.
\end{proof}

\begin{corollary}[Quaternion fixed-excess fixed-\(k\) column-Gram lower side]\label{cor:quat-rect-kth-gram-lower}
Under the hypotheses of Theorem~\ref{thm:quat-fixed-excess-fixed-k-lower}, let
\[
  \Lambda_{\HH,m}^{(r)}[k]
  =
  \min_{\substack{T\subset[N_m]\\ |T|=m+r}}
  \lambda_{(k)}\!\left((m+r)^{-1}A_{m,T}^*A_{m,T}\right),
\]
where \(\lambda_{(k)}\) is the \(k\)-th smallest quaternionic eigenvalue of the
Hermitian column Gram matrix.  If \(N_m/m\to\gamma>1\), then, for every
\(\varepsilon>0\),
\[
  \PP\left\{\Lambda_{\HH,m}^{(r)}[k]\le
  \exp\!\left[-\left(\frac{h(\gamma)}{2k(k+r)}+\varepsilon\right)m\right]\right\}
  \to0,
\]
and, on any common coupling,
\[
  \liminf_{m\to\infty}\frac1m\log \Lambda_{\HH,m}^{(r)}[k]
  \ge -\frac{h(\gamma)}{2k(k+r)}
  \qquad\hbox{a.s.}
\]
\end{corollary}

\begin{proof}
Pathwise,
\[
  \Lambda_{\HH,m}^{(r)}[k]
  =(m+r)^{-1}\bigl(M_{\HH,m}^{(r)}[k]\bigr)^2,
\]
and \(m^{-1}\log(m+r)\to0\).  Apply
Theorem~\ref{thm:quat-fixed-excess-fixed-k-lower}.
\end{proof}

\begin{corollary}[Random-prefix fixed-excess fixed-\(k\) lower sides]\label{cor:random-prefix-fixed-excess-fixed-k-lower}
Let \(\F\in\{\R,\C,\HH\}\), with \(d_\F=1,2,4\), respectively.  For each
\(m\), let \((R_{im})_{i\ge1}\) be an infinite sequence of iid standard
\(\F^m\)-Gaussian rows.  Fix \(r\ge0\) and \(k\ge1\).  For an integer
\(n\ge m+r\), put
\[
  M_{\F,m}^{(r)}[k](n)
  =
  \min_{\substack{T\subset[n]\\ |T|=m+r}}
  \sigma_{m-k+1}((R_{im})_{i\in T}),
\]
and
\[
  \Lambda_{\F,m}^{(r)}[k](n)
  =
  \min_{\substack{T\subset[n]\\ |T|=m+r}}
  \lambda_{(k)}\!\left((m+r)^{-1}(R_{im})_{i\in T}^*
  (R_{im})_{i\in T}\right).
\]
Let \(N_m\) be integer-valued, possibly dependent on the row array, with
\[
  \frac{N_m}{m}\xrightarrow{\PP}\gamma>1 .
\]
Extend both statistics by \(+\infty\) on the event \(N_m<m+r\).  This event has
probability tending to zero; the convention only makes the monotonicity
comparison literal.
Then, for every \(\varepsilon>0\),
\[
  \PP\left\{M_{\F,m}^{(r)}[k](N_m)\le
  \exp\!\left[-\left(\frac{h(\gamma)}{d_\F k(k+r)}
  +\varepsilon\right)m\right]\right\}\to0
\]
and
\[
  \PP\left\{\Lambda_{\F,m}^{(r)}[k](N_m)\le
  \exp\!\left[-\left(\frac{2h(\gamma)}{d_\F k(k+r)}
  +\varepsilon\right)m\right]\right\}\to0.
\]
On a common coupling, if \(N_m/m\to\gamma\) almost surely, then
\[
  \liminf_{m\to\infty}\frac1m\log M_{\F,m}^{(r)}[k](N_m)
  \ge -\frac{h(\gamma)}{d_\F k(k+r)}
\]
and
\[
  \liminf_{m\to\infty}\frac1m\log \Lambda_{\F,m}^{(r)}[k](N_m)
  \ge -\frac{2h(\gamma)}{d_\F k(k+r)}
  \qquad\hbox{a.s.}
\]
\end{corollary}

\begin{proof}
The map \(n\mapsto M_{\F,m}^{(r)}[k](n)\) is nonincreasing, because increasing
the prefix row count enlarges the tested family.  Fix \(\varepsilon>0\), and
choose \(\gamma_+>\gamma\) so close that
\[
  \frac{h(\gamma_+)}{d_\F k(k+r)}
  <
  \frac{h(\gamma)}{d_\F k(k+r)}+\frac{\varepsilon}{2}.
\]
With probability tending to one,
\[
  N_m\le \lceil\gamma_+m\rceil .
\]
On this event, including the convention \(M=+\infty\) when \(N_m<m+r\),
\[
  M_{\F,m}^{(r)}[k](N_m)
  \ge M_{\F,m}^{(r)}[k](\lceil\gamma_+m\rceil).
\]
Apply the deterministic-count fixed-aspect lower side
Corollary~\ref{rect-kth:fixed-aspect-lower} for \(\F=\R,\C\), and
Theorem~\ref{thm:quat-fixed-excess-fixed-k-lower} for \(\F=\HH\), at the
deterministic aspect \(\gamma_+\).  At the threshold
\[
  \exp\left[-\left(\frac{h(\gamma)}{d_\F k(k+r)}+\varepsilon\right)m\right],
\]
this deterministic lower side applies with at least \(\varepsilon/2\) slack.
No independence between \(N_m\) and the rows is used.

The column-Gram statement is pathwise:
\[
  \Lambda_{\F,m}^{(r)}[k](N_m)
  =(m+r)^{-1}\bigl(M_{\F,m}^{(r)}[k](N_m)\bigr)^2,
\]
on the event \(N_m\ge m+r\), while both sides are \(+\infty\) under the
extension on the complement; \(m^{-1}\log(m+r)\to0\).  For the almost-sure
lower envelopes, fix a countable decreasing sequence \(\gamma_j\downarrow\gamma\).
Almost surely, eventually \(N_m\le\lceil\gamma_jm\rceil\) for each \(j\);
combine the preceding pathwise comparison with the deterministic-count
almost-sure lower envelope at \(\gamma_j\), then let \(j\to\infty\).

This is only the lower-side monotonicity transfer.  In the overlapping square
case \(r=0,k=1\), the exact random-prefix laws are the earlier square results;
the present corollary adds fixed-excess fixed-\(k\) lower sides and normalized
column-Gram lower sides only.  It does not assert an upper side, exact
fixed-excess law, value-dependent row-selection theorem, row-Gram statement for
\(r>0\), beta interpolation, or proportional Wishart-edge law.
\end{proof}

\begin{corollary}[Random-prefix diagonal fixed-excess fixed-\(k\) lower sides]\label{cor:random-prefix-diagonal-fixed-excess-fixed-k}
Let \(\F\in\{\R,\C,\HH\}\), with \(d_\F=1,2,4\), respectively.  For each
\(m\), let \((R_{im})_{i\ge1}\) be an infinite sequence of iid standard
\(\F^m\)-Gaussian rows.  Let \(N_m\) be integer-valued, possibly dependent on
the row array, with
\[
  \frac{N_m}{m}\xrightarrow{\PP}\gamma>1 .
\]
Fix \(r\ge0\) and \(k\ge1\).  Let \(d_{im}>0\), \(i\ge1\), and
\(e_{jm}>0\), \(1\le j\le m\), be finite real random weights, possibly
dependent on the rows and on \(N_m\).  Assume
\[
  \Delta_m^{\rm rpdiag}
  =
  \max_{1\le i\le N_m}|\log d_{im}|
  +\max_{1\le j\le m}|\log e_{jm}|
  =o_\PP(m),
\]
with the convention that an absent row or column factor is replaced by weights
identically equal to \(1\).  For \(N_m\ge m+r\) and \(T\subset[N_m]\),
\(|T|=m+r\), set
\[
  B_{m,T}=\operatorname{diag}(d_{im}:i\in T)\,(R_{im})_{i\in T}\,
  \operatorname{diag}(e_{1m},\ldots,e_{mm})
\]
and define
\[
  M_{\rm rpdiag,\F}^{(r)}[k]
  =
  \min_{\substack{T\subset[N_m]\\ |T|=m+r}}
  \sigma_{m-k+1}(B_{m,T}),
\]
\[
  \Lambda_{\rm rpdiag,\F}^{(r)}[k]
  =
  \min_{\substack{T\subset[N_m]\\ |T|=m+r}}
  \lambda_{(k)}\!\left((m+r)^{-1}B_{m,T}^*B_{m,T}\right).
\]
Extend both statistics by \(+\infty\) on \(\{N_m<m+r\}\).  Then, for every
\(\varepsilon>0\),
\[
  \PP\left\{M_{\rm rpdiag,\F}^{(r)}[k]\le
  \exp\!\left[-\left(\frac{h(\gamma)}{d_\F k(k+r)}
    +\varepsilon\right)m\right]\right\}\to0
\]
and
\[
  \PP\left\{\Lambda_{\rm rpdiag,\F}^{(r)}[k]\le
  \exp\!\left[-\left(\frac{2h(\gamma)}{d_\F k(k+r)}
    +\varepsilon\right)m\right]\right\}\to0.
\]
On a common coupling, if \(N_m/m\to\gamma\) almost surely and
\(\Delta_m^{\rm rpdiag}=o(m)\) almost surely, then
\[
  \liminf_{m\to\infty}\frac1m\log M_{\rm rpdiag,\F}^{(r)}[k]
  \ge -\frac{h(\gamma)}{d_\F k(k+r)}
\]
and
\[
  \liminf_{m\to\infty}\frac1m\log \Lambda_{\rm rpdiag,\F}^{(r)}[k]
  \ge -\frac{2h(\gamma)}{d_\F k(k+r)}
  \qquad\hbox{a.s.}
\]
The same conclusions apply to row-normalized random-prefix Gaussian rows and
to full-prefix column-normalized Gaussian columns, where the column
normalization is computed from all rows in the realized prefix \([N_m]\).  For
independent thinning, this is read only after the retained iid rows are
reindexed; value-dependent thinning or row selection is not covered.
\end{corollary}

\begin{proof}
On the feasible event \(\{N_m\ge m+r\}\), for every selected row set \(T\),
\[
  \sigma_{m-k+1}(B_{m,T})
  \ge
  e^{-\Delta_m^{\rm rpdiag}}
  \sigma_{m-k+1}\!\left((R_{im})_{i\in T}\right).
\]
Taking minima over \(T\) gives the same lower sandwich for the random-prefix
singular-value statistics, and the normalized column-Gram statistic has the
squared sandwich.  Hence, with
\(c=h(\gamma)/(d_\F k(k+r))\),
\[
  \PP\{M_{\rm rpdiag,\F}^{(r)}[k]\le e^{-(c+\varepsilon)m}\}
  \le
  \PP\{\Delta_m^{\rm rpdiag}>\varepsilon m/2\}
  +\PP\{M_{\F,m}^{(r)}[k](N_m)\le e^{-(c+\varepsilon/2)m}\}.
\]
The first term tends to zero by the diagonal-extrema hypothesis, and the second
by Corollary~\ref{cor:random-prefix-fixed-excess-fixed-k-lower}.  The Gram
claim follows in the same way, using the lower sandwich
\(\Lambda_{\rm rpdiag,\F}^{(r)}[k]\ge
e^{-2\Delta_m^{\rm rpdiag}}\Lambda_{\F,m}^{(r)}[k](N_m)\) on the feasible
event and splitting the slack accordingly.  The almost-sure lower envelopes use
the same pathwise inequalities together with the a.s. part of
Corollary~\ref{cor:random-prefix-fixed-excess-fixed-k-lower}.

For row normalization, on the high-probability bracket
\(N_m\le\lceil\gamma_+m\rceil\), the row-norm extrema over the realized prefix
are bounded by the extrema over the deterministic prefix
\(\lceil\gamma_+m\rceil\).  For column normalization, on
\(\lfloor\gamma_-m\rfloor\le N_m\le\lceil\gamma_+m\rceil\), each full-prefix
column norm is bracketed between the corresponding deterministic-prefix column
norms.  The same Gaussian norm upper-tail and ball-volume lower-tail estimates
used in the fixed-count normalization corollaries give
\(\Delta_m^{\rm rpdiag}=o_\PP(m)\) and, under common a.s. count convergence,
\(\Delta_m^{\rm rpdiag}=o(m)\) almost surely.

This is only a random-prefix diagonal lower-side transfer.  In the overlapping
square case \(r=0,k=1\), the exact random-prefix square laws above are stronger.
No upper side, exact fixed-excess law, value-dependent row-selection theorem,
row-Gram result for \(r>0\), selected-submatrix column normalization,
non-diagonal covariance statement, beta interpolation, proportional Wishart
edge, or supercritical phase-retrieval consequence is implied.
\end{proof}

\begin{corollary}[Random-prefix all-column fixed-excess fixed-\(k\) lower sides]\label{cor:random-prefix-right-factor-fixed-excess-fixed-k-lower}
Let \(\F\in\{\R,\C,\HH\}\), with \(d_\F=1,2,4\), respectively.  For each
\(m\), let \((R_{im})_{i\ge1}\) be an infinite sequence of iid standard
\(\F^m\)-Gaussian rows.  Let \(N_m\) be integer-valued, possibly dependent on
the row array, with
\[
  \frac{N_m}{m}\xrightarrow{\PP}\gamma>1 .
\]
Fix \(r\ge0\) and \(k\ge1\).  Let \(d_{im}>0\), \(i\ge1\), be finite real
random row weights, and let \(S_m\in\F^{m\times m}\) be invertible.  The row
weights and \(S_m\) may be random and dependent on \(N_m\) and on the row array.
Assume
\[
  \Delta_m^{\rm rpcond}
  =
  \max_{1\le i\le N_m}|\log d_{im}|
  +\bigl|\log\smin(S_m)\bigr|+\bigl|\log\|S_m\|\bigr|
  =o_\PP(m),
\]
omitting the row-weight term when row weights are absent.  For \(N_m\ge m+r\)
and \(T\subset[N_m]\), \(|T|=m+r\), set
\[
  B_{m,T}=\operatorname{diag}(d_{im}:i\in T)\,(R_{im})_{i\in T}S_m
\]
and define
\[
  M_{\rm rpcond,\F}^{(r)}[k]
  =
  \min_{\substack{T\subset[N_m]\\ |T|=m+r}}\sigma_{m-k+1}(B_{m,T}),
\]
\[
  \Lambda_{\rm rpcond,\F}^{(r)}[k]
  =
  \min_{\substack{T\subset[N_m]\\ |T|=m+r}}
  \lambda_{(k)}\!\left((m+r)^{-1}B_{m,T}^*B_{m,T}\right).
\]
Extend both statistics by \(+\infty\) on \(\{N_m<m+r\}\).  Then, for every
\(\varepsilon>0\),
\[
  \PP\left\{M_{\rm rpcond,\F}^{(r)}[k]\le
  \exp\!\left[-\left(\frac{h(\gamma)}{d_\F k(k+r)}
    +\varepsilon\right)m\right]\right\}\to0
\]
and
\[
  \PP\left\{\Lambda_{\rm rpcond,\F}^{(r)}[k]\le
  \exp\!\left[-\left(\frac{2h(\gamma)}{d_\F k(k+r)}
    +\varepsilon\right)m\right]\right\}\to0.
\]
On a common coupling, if \(N_m/m\to\gamma\) almost surely and
\(\Delta_m^{\rm rpcond}=o(m)\) almost surely, then
\[
  \liminf_{m\to\infty}\frac1m\log M_{\rm rpcond,\F}^{(r)}[k]
  \ge -\frac{h(\gamma)}{d_\F k(k+r)}
\]
and
\[
  \liminf_{m\to\infty}\frac1m\log \Lambda_{\rm rpcond,\F}^{(r)}[k]
  \ge -\frac{2h(\gamma)}{d_\F k(k+r)}
  \qquad\hbox{a.s.}
\]
Positive column diagonals are included by taking \(S_m\) diagonal.  Deterministic
or random full-column covariance and whitening factors are included only under
the absolute singular-scale hypothesis above.  Row-normalized random-prefix
Gaussian rows and full-prefix column-normalized Gaussian columns may be combined
with this all-column right factor when their realized-prefix logarithmic extrema
are included in \(\Delta_m^{\rm rpcond}\).
\end{corollary}

\begin{proof}
On the feasible event \(\{N_m\ge m+r\}\), for every selected row set \(T\),
\[
  \sigma_{m-k+1}(B_{m,T})
  \ge
  e^{-\Delta_m^{\rm rpcond}}
  \sigma_{m-k+1}\!\left((R_{im})_{i\in T}\right).
\]
For \(\F=\HH\), the positive row weights are real scalars and hence central,
while \(S_m\) is one global quaternionic right factor; the same min--max
singular-value sandwich applies.  Taking minima over \(T\) gives the lower
sandwich for the random-prefix singular-value statistics.  With
\(c=h(\gamma)/(d_\F k(k+r))\),
\[
  \PP\{M_{\rm rpcond,\F}^{(r)}[k]\le e^{-(c+\varepsilon)m}\}
  \le
  \PP\{\Delta_m^{\rm rpcond}>\varepsilon m/2\}
  +\PP\{M_{\F,m}^{(r)}[k](N_m)\le e^{-(c+\varepsilon/2)m}\},
\]
and both terms tend to zero by the conditioning hypothesis and
Corollary~\ref{cor:random-prefix-fixed-excess-fixed-k-lower}.  No independence
among the count, row array, row weights, and \(S_m\) is used.

For the normalized column Gram, the identity
\[
  \lambda_{(k)}\!\left((m+r)^{-1}B_{m,T}^*B_{m,T}\right)
  =(m+r)^{-1}\sigma_{m-k+1}(B_{m,T})^2
\]
gives
\[
  \Lambda_{\rm rpcond,\F}^{(r)}[k]\ge
  e^{-2\Delta_m^{\rm rpcond}}\Lambda_{\F,m}^{(r)}[k](N_m)
\]
on the feasible event.  Hence
\[
  \PP\{\Lambda_{\rm rpcond,\F}^{(r)}[k]\le e^{-(2c+\varepsilon)m}\}
  \le
  \PP\{\Delta_m^{\rm rpcond}>\varepsilon m/4\}
  +\PP\{\Lambda_{\F,m}^{(r)}[k](N_m)\le e^{-(2c+\varepsilon/2)m}\},
\]
which tends to zero by the same two inputs.  The almost-sure lower envelopes
follow by intersecting the a.s. lower-envelope event from
Corollary~\ref{cor:random-prefix-fixed-excess-fixed-k-lower} with
\(\Delta_m^{\rm rpcond}=o(m)\).

This is only a random-prefix row-submatrix lower-side transfer under one global
all-column right factor.  A subexponential condition number alone is
insufficient, since a scalar factor \(e^{cm}I\) shifts the exponential scale.
The statement does not assert an upper side, fixed-excess exact law,
value-dependent row-selection theorem, selected-column covariance transfer,
selected-submatrix column normalization, row-Gram statement for \(r>0\), beta
interpolation, proportional Wishart-edge law, or supercritical phase-retrieval
consequence.
\end{proof}

\begin{conjecture}[Quaternionic fixed-excess hard edge]\label{conj:quat-fixed-excess}
Let \(A_m\in\HH^{N_m\times m}\) have iid standard quaternion Gaussian entries
and \(N_m/m\to\gamma>1\).  For each fixed integer \(r\ge1\), one expects
\[
  \frac{1}{m}\log\min_{\substack{T\subset[N_m]\\ |T|=m+r}}\smin(A_{m,T})
  \xrightarrow{\PP}-\frac{h(\gamma)}{4(r+1)}.
\]
The lower side is the \(k=1\) case of
Theorem~\ref{thm:quat-fixed-excess-fixed-k-lower}; the matching upper side still
requires a quaternionic fixed-excess pair-correlation theorem.
\end{conjecture}

\begin{remark}
The quaternion results above use quaternionic singular values, not the repeated
realification indices.  They are lower-side extensions except for the square
least-singular-value theorem and its positive-real diagonal/normalization
transfer and square all-column right-factor transfer, whose upper sides are
proved separately.  They do not assert a quaternionic version of the
spherical-radial transfer theorem, whose statement in this paper is restricted
to \(\F\in\{\R,\C\}\), and they do not promote fixed-excess quaternion exact
laws, selected-column non-diagonal covariance transfers, proportional-excess
Wishart-edge laws, beta-interpolation statements, or row-Gram statements for
\(r>0\).
\end{remark}

\subsection{Deterministic row and column conditioning}

\begin{proposition}[Subexponential deterministic conditioning transfer]\label{prop:cov-transfer}
Let \(G_m\in\F^{N_m\times m}\), \(N_m/m\to\gamma>1\), have iid standard
Gaussian entries.  Let \(D_m=\operatorname{diag}(d_{1m},\dots,d_{N_m,m})\) be
deterministic with positive diagonal entries, let \(\Sigma_m\) be deterministic
positive definite, and put
\[
  A_m=D_mG_m\Sigma_m^{1/2}.
\]
Assume
\[
  \frac1m\log\!\left(\min_i d_{im}\sqrt{\lambda_{\min}(\Sigma_m)}\right)\to0,
  \qquad
  \frac1m\log\!\left(\max_i d_{im}\sqrt{\lambda_{\max}(\Sigma_m)}\right)\to0.
\]
Then
\[
  \frac{1}{m}\log\min_{|T|=m}\smin(A_{m,T})
  -
  \frac{1}{m}\log\min_{|T|=m}\smin(G_{m,T})
  \to0
\]
deterministically.  Consequently
\[
  \frac{1}{m}\log\min_{|T|=m}\smin(A_{m,T})
  \xrightarrow{\PP}-\frac{h(\gamma)}{d_\F}.
\]
\end{proposition}

\begin{proof}
For every row set \(T\), write \(D_T\) for the diagonal submatrix of \(D_m\) on
the selected rows.  Then
\[
  A_{m,T}=D_TG_{m,T}\Sigma_m^{1/2},
\]
and therefore
\[
  \left(\min_i d_{im}\right)\sqrt{\lambda_{\min}(\Sigma_m)}\,\smin(G_{m,T})
  \le \smin(A_{m,T})
  \le
  \left(\max_i d_{im}\right)\sqrt{\lambda_{\max}(\Sigma_m)}\,\smin(G_{m,T}).
\]
Taking minima over \(T\), logs, and dividing by \(m\) gives the deterministic
comparison.  The limit then follows from the Gaussian source theorem.
\end{proof}

\begin{corollary}[Uniformly elliptic row-scaled Gaussian frames]\label{cor:elliptic-gaussian}
If, for constants \(0<c<C<\infty\),
\[
  c\le d_{im}\le C,\qquad
  c\le\lambda_{\min}(\Sigma_m)\le\lambda_{\max}(\Sigma_m)\le C
\]
for all \(i,m\), then the extremal row-submatrix base is the Gaussian base
\(h(\gamma)/d_\F\).  This includes Gaussian frames with common covariance and
bounded deterministic row weights.
\end{corollary}

\begin{proof}
The uniform ellipticity assumptions imply the two logarithmic conditioning
hypotheses in Proposition~\ref{prop:cov-transfer}.
\end{proof}

\begin{proposition}[Random row-weight transfer]\label{prop:random-row-weight}
Let \(A_m\in\F^{N_m\times n_m}\) be arbitrary random matrices and let
\(\mathcal T_m\) be any family of row sets for which the selected least
singular values are defined.  Put
\[
  M_m(A)=\min_{T\in\mathcal T_m}\smin((A_m)_T).
\]
Let \(D_m^{\rm row}=\operatorname{diag}(w_{1m},\ldots,w_{N_m,m})\) have positive
random diagonal entries, possibly dependent on \(A_m\), and set
\(B_m=D_m^{\rm row}A_m\).  For
any scale \(a_m\to\infty\), define
\[
  \Delta_m=\max_{1\le i\le N_m}|\log w_{im}|.
\]
On the event \(M_m(A)>0\),
\[
  \left|\frac1{a_m}\log M_m(B)-\frac1{a_m}\log M_m(A)\right|
  \le \frac{\Delta_m}{a_m}.
\]
Consequently, if \(\PP\{M_m(A)>0\}\to1\),
\(\Delta_m=o_\PP(a_m)\), and
\(a_m^{-1}\log M_m(A)\to L\) in probability, then
\[
  \frac1{a_m}\log M_m(B)\xrightarrow{\PP} L.
\]
\end{proposition}

\begin{proof}
For every row set \(T\),
\[
  (B_m)_T=(D_m^{\rm row})_T(A_m)_T.
\]
Since \(e^{-\Delta_m}\le w_{im}\le e^{\Delta_m}\) for every row,
\[
  e^{-\Delta_m}\smin((A_m)_T)
  \le \smin((B_m)_T)
  \le e^{\Delta_m}\smin((A_m)_T).
\]
The constants are independent of \(T\).  Taking minima over \(T\), logs, and
division by \(a_m\) gives the displayed comparison.
\end{proof}

\begin{proposition}[Random column-weight transfer]\label{prop:random-column-weight}
Let \(A_m\in\F^{N_m\times n_m}\) be arbitrary random matrices.  Let
\(\mathcal I_m\) be any nonempty family of row-column pairs \((R,C)\) for which the
least singular values of \((A_m)_{R,C}\) are defined, and put
\[
  M_m(A)=\min_{(R,C)\in\mathcal I_m}\smin((A_m)_{R,C}).
\]
Let \(E_m=\operatorname{diag}(e_{1m},\ldots,e_{n_m,m})\) have strictly positive
finite random diagonal entries, possibly dependent on \(A_m\), and set
\(B_m=A_mE_m\).  For any scale \(a_m\to\infty\), define
\[
  \Delta_m^{\rm col}=\max_{1\le j\le n_m}|\log e_{jm}|.
\]
On the event \(M_m(A)>0\),
\[
  \left|\frac1{a_m}\log M_m(B)-\frac1{a_m}\log M_m(A)\right|
  \le \frac{\Delta_m^{\rm col}}{a_m}.
\]
Consequently, if \(\PP\{M_m(A)>0\}\to1\),
\(\Delta_m^{\rm col}=o_\PP(a_m)\), and
\(a_m^{-1}\log M_m(A)\to L\) in probability, then
\[
  \frac1{a_m}\log M_m(B)\xrightarrow{\PP} L.
\]
\end{proposition}

\begin{proof}
For every selected pair \((R,C)\),
\[
  (B_m)_{R,C}=(A_m)_{R,C}(E_m)_C ,
\]
where \((E_m)_C\) is the positive diagonal submatrix of \(E_m\) on the selected
columns.  Since
\[
  e^{-\Delta_m^{\rm col}}\le e_{jm}\le e^{\Delta_m^{\rm col}}
  \qquad(1\le j\le n_m),
\]
the singular-value inequalities for a right diagonal factor give
\[
  e^{-\Delta_m^{\rm col}}\smin((A_m)_{R,C})
  \le \smin((B_m)_{R,C})
  \le e^{\Delta_m^{\rm col}}\smin((A_m)_{R,C}).
\]
The constants are uniform over \((R,C)\).  Taking minima, logs, and division by
\(a_m\) proves the comparison.
\end{proof}

\begin{corollary}[Gaussian row-normalized frames]\label{cor:row-normalized-gaussian}
Let \(G_m\in\F^{N_m\times m}\) have iid standard Gaussian entries, with
\(N_m/m\to\gamma>1\).  Let \(U_m\) be the row-normalized matrix
\[
  (U_m)_i=\frac{(G_m)_i}{\|(G_m)_i\|_2}.
\]
Then
\[
  \frac1m\log\min_{|T|=m}\smin((U_m)_T)
  \xrightarrow{\PP}-\frac{h(\gamma)}{d_\F}.
\]
On any common coupling over \(m\), the lower envelope also satisfies
\[
  \liminf_{m\to\infty}\frac1m\log\min_{|T|=m}\smin((U_m)_T)
  \ge -\frac{h(\gamma)}{d_\F}
  \qquad\hbox{almost surely}.
\]
\end{corollary}

\begin{proof}
Apply Proposition~\ref{prop:random-row-weight} with
\[
  w_{im}=\|(G_m)_i\|_2^{-1},\qquad a_m=m.
\]
The Gaussian rows are nonzero almost surely, and all square Gaussian row
submatrices are nonsingular almost surely, so the positive-weight and
\(M_m(G)>0\) hypotheses of the transfer proposition hold with probability one.
It remains to check
\[
  \max_i|\log\|(G_m)_i\|_2|=o_\PP(m).
\]
For every fixed \(\delta>0\), standard Gaussian norm tails give
\[
  \PP\left\{\max_{i\le N_m}\|(G_m)_i\|_2>e^{\delta m}\right\}\to0.
\]
For the lower tail, a ball-volume estimate in real dimension \(d_\F m\) gives
\[
  \PP\left\{\min_{i\le N_m}\|(G_m)_i\|_2<e^{-\delta m}\right\}
  \le \exp\{-d_\F\delta m^2+O(m\log m)\}\to0.
\]
Thus the random row weights have subexponential extrema, and the Gaussian
source theorem transfers to \(U_m\).

For the almost-sure lower envelope, only the upper row-norm extreme is needed:
\[
  \min_{|T|=m}\smin((U_m)_T)
  \ge
  \frac{\min_{|T|=m}\smin((G_m)_T)}
       {\max_{i\le N_m}\|(G_m)_i\|_2}.
\]
Since \(N_m=O(m)\), Markov's inequality gives
\[
  \sum_m
  \PP\left\{\max_{i\le N_m}\|(G_m)_i\|_2>e^{\delta m}\right\}
  \le
  \sum_m N_m\,\EE\|(G_m)_1\|_2^2 e^{-2\delta m}<\infty .
\]
Corollary~\ref{cor:as-square-lower-envelope}, Borel--Cantelli for this
row-norm bound, and then \(\delta\downarrow0\) give the almost-sure lower
envelope.
\end{proof}

\begin{remark}[Normalized non-Gaussian rows]
For lower-side estimates, one only needs the upper row-norm extreme.  If
\(B_i=A_i/\|A_i\|_2\), the rows are nonzero with high probability, and
\[
  \max_{i\le N_m}\|A_i\|_2=e^{o_\PP(m)},
\]
then \(M_m(B)\ge e^{-o_\PP(m)}M_m(A)\).  Hence any lower-side exponential
bound already proved for \(A_m\) transfers to the row-normalized ensemble.
Bounded scalar density alone does not imply this row-tail condition for
pathologically heavy-tailed laws; it should be assumed, or replaced by a
sufficient tail hypothesis such as
\[
  N_m m\,\PP\{|X_m|>e^{\delta m}/\sqrt m\}\to0
  \qquad(\delta>0).
\]
The normalized rows are sphere-supported and their coordinates are dependent,
so the bounded-density entry theorem cannot be applied directly to \(B_m\).
\end{remark}

\begin{corollary}[Gaussian column-normalized matrices]\label{cor:column-normalized-gaussian}
Let \(G_m\in\F^{N_m\times m}\) have iid standard Gaussian entries, with
\(N_m/m\to\gamma>1\).  Let \(V_m\) be the column-normalized matrix
\[
  (V_m)^{(j)}=\frac{(G_m)^{(j)}}{\|(G_m)^{(j)}\|_2},
  \qquad 1\le j\le m,
\]
where \((G_m)^{(j)}\) denotes the \(j\)-th column.  Then
\[
  \frac1m\log\min_{|T|=m}\smin((V_m)_T)
  \xrightarrow{\PP}-\frac{h(\gamma)}{d_\F}.
\]
On any common coupling over \(m\), the lower envelope also satisfies
\[
  \liminf_{m\to\infty}\frac1m\log\min_{|T|=m}\smin((V_m)_T)
  \ge -\frac{h(\gamma)}{d_\F}
  \qquad\hbox{almost surely}.
\]
\end{corollary}

\begin{proof}
Write \(V_m=G_mE_m\), where
\[
  E_m=\operatorname{diag}\bigl(\|(G_m)^{(1)}\|_2^{-1},\ldots,
  \|(G_m)^{(m)}\|_2^{-1}\bigr).
\]
The Gaussian columns are nonzero almost surely.  We apply
Proposition~\ref{prop:random-column-weight} to the row-column family
\((T,[m])\), \(|T|=m\); Gaussian square row submatrices are nonsingular almost
surely, so the required positivity event has probability one.
By Proposition~\ref{prop:random-column-weight}, it is enough for convergence in
probability to prove
\[
  \max_{1\le j\le m}\left|\log\|(G_m)^{(j)}\|_2\right|=o_\PP(m).
\]
For the almost-sure lower envelope, the same estimates will be summable.  For
the upper tail, Markov's inequality gives, for every fixed \(\delta>0\),
\[
  \PP\{\|(G_m)^{(j)}\|_2>e^{\delta m}\}
  \le N_m e^{-2\delta m},
\]
and the union over \(j\le m\) is summable in \(m\).  For the lower tail, the
column lives in real dimension \(d_\F N_m\), so a ball-volume estimate gives
\[
  \PP\{\|(G_m)^{(j)}\|_2<e^{-\delta m}\}
  \le \exp\{-d_\F\delta N_m m+O(N_m\log N_m)\}.
\]
Since \(N_m/m\to\gamma>1\), the union over \(j\le m\) is again summable.  Thus
the column weights have subexponential extrema in probability and almost
surely.  The Gaussian source theorem gives the convergence in probability, and
Corollary~\ref{cor:as-square-lower-envelope} plus the pathwise sandwich gives
the almost-sure lower envelope.
\end{proof}

\begin{remark}[Product variance profiles]
For real Gaussian entries \(A_{ij}=d_i g_{ij}c_j\), this is a bounded product
variance profile.  It is not a theorem for arbitrary variance profiles
\(\sigma_{ij}g_{ij}\): row types may have different entropy and small-ball
prefactors, so the sharp base should be profile-dependent.
\end{remark}

\begin{corollary}[Fixed-excess transfer]\label{cor:fixed-excess-transfer}
Under the hypotheses of Proposition~\ref{prop:cov-transfer}, fix \(r\ge0\) and
define
\[
  M_A^{(r)}=\min_{|T|=m+r}\smin(A_{m,T}),
  \qquad
  M_G^{(r)}=\min_{|T|=m+r}\smin(G_{m,T}).
\]
Then
\[
  \frac1m\log M_A^{(r)}-\frac1m\log M_G^{(r)}\to0
\]
deterministically.  Consequently the fixed-excess lower side of
Theorem~\ref{rect:lower} transfers to \(A_m\): for every \(\varepsilon>0\),
\[
  \PP\left\{M_A^{(r)}\le
     \exp\!\left[-\left(\frac{h(\gamma)}{d_\F(r+1)}+\varepsilon\right)m\right]\right\}
  \to0.
\]
\end{corollary}

\begin{proof}
The singular-value sandwich in the proof of Proposition~\ref{prop:cov-transfer}
does not use \(|T|=m\).  It applies uniformly to every \((m+r)\)-row set \(T\).
Taking minima, logs, and dividing by \(m\) gives the deterministic comparison,
and Theorem~\ref{rect:lower} gives the lower side.
\end{proof}

\begin{corollary}[Fixed-excess fixed-\(k\) transfer]\label{cor:fixed-excess-fixed-k-transfer}
Under the hypotheses of Proposition~\ref{prop:cov-transfer}, with
\(\F\in\{\R,\C\}\), fix integers \(r\ge0\) and \(k\ge1\).  For \(m\ge k\), put
\[
  M_A^{(r)}[k]=
  \min_{\substack{T\subset[N_m]\\ |T|=m+r}}\sigma_{m-k+1}(A_{m,T}),
  \qquad
  M_G^{(r)}[k]=
  \min_{\substack{T\subset[N_m]\\ |T|=m+r}}\sigma_{m-k+1}(G_{m,T}).
\]
Then
\[
  \frac1m\log M_A^{(r)}[k]-\frac1m\log M_G^{(r)}[k]\to0
\]
deterministically.  Consequently the fixed-excess fixed-\(k\) lower side of
Corollary~\ref{rect-kth:fixed-aspect-lower} transfers to \(A_m\): for every
\(\varepsilon>0\),
\[
  \PP\left\{M_A^{(r)}[k]\le
  \exp\!\left[-\left(\frac{h(\gamma)}{d_\F k(k+r)}+\varepsilon\right)m\right]\right\}
  \to0.
\]
The normalized column-Gram statistic
\[
  \Lambda_A^{(r)}[k]
  =
  \min_{\substack{T\subset[N_m]\\ |T|=m+r}}
  \lambda_{(k)}\!\left((m+r)^{-1}A_{m,T}^*A_{m,T}\right)
\]
has the same deterministic transfer from the corresponding statistic
\(\Lambda_G^{(r)}[k]\) defined with \(G_m\) in place of \(A_m\), and hence has
the lower-side base \(2h(\gamma)/(d_\F k(k+r))\) under the same hypotheses.
On any common coupling over \(m\), the lower envelopes also satisfy
\[
  \liminf_{m\to\infty}\frac1m\log M_A^{(r)}[k]
  \ge -\frac{h(\gamma)}{d_\F k(k+r)}
  \qquad\hbox{almost surely},
\]
and
\[
  \liminf_{m\to\infty}\frac1m\log \Lambda_A^{(r)}[k]
  \ge -\frac{2h(\gamma)}{d_\F k(k+r)}
  \qquad\hbox{almost surely}.
\]
\end{corollary}

\begin{proof}
For every row set \(T\), the min--max singular-value inequalities applied to
\[
  A_{m,T}=D_TG_{m,T}\Sigma_m^{1/2}
\]
give
\[
\begin{aligned}
  \left(\min_i d_{im}\right)\sqrt{\lambda_{\min}(\Sigma_m)}\,
  \sigma_{m-k+1}(G_{m,T})
  &\le \sigma_{m-k+1}(A_{m,T})\\
  &\le
  \left(\max_i d_{im}\right)\sqrt{\lambda_{\max}(\Sigma_m)}\,
  \sigma_{m-k+1}(G_{m,T}).
\end{aligned}
\]
The constants are independent of \(T\).  Taking minima, logs, and dividing by
\(m\) gives the deterministic comparison, and
Corollary~\ref{rect-kth:fixed-aspect-lower} gives the lower side.  Squaring the
same comparison and multiplying by \((m+r)^{-1}\) gives the normalized
column-Gram transfer, using Corollary~\ref{rect-kth:gram-lower}.  On a common
coupling over \(m\), Corollary~\ref{cor:as-geometric-lower-envelopes} gives the
Gaussian lower envelopes for \(M_G^{(r)}[k]\) and \(\Lambda_G^{(r)}[k]\).  The
deterministic comparison above and its squared Gram version have \(o(1)\)
logarithmic error after division by \(m\), so the two liminf bounds transfer to
\(A_m\).  Absolute extreme eigenvalue control is used here; a subexponential
condition number alone would not prevent a global exponential rescaling.
\end{proof}

\begin{remark}[Random diagonal weights at fixed excess]\label{rem:random-diagonal-fixed-k}
The same min--max argument used in
Propositions~\ref{prop:random-row-weight} and~\ref{prop:random-column-weight}
applies to the fixed-excess fixed-\(k\) statistics above.  Let
\[
  \widetilde A_m=D_m^{\rm row}A_mE_m,
\]
where the diagonal entries are strictly positive and finite, random and
possibly dependent on \(A_m\), and set
\[
  \Delta_m^{\rm diag}
  =
  \max_{i\le N_m}|\log (D_m^{\rm row})_{ii}|
  +
  \max_{j\le m}|\log (E_m)_{jj}|,
\]
omitting the absent term if only one side is weighted.  On the event
\(M_A^{(r)}[k]>0\),
\[
  \left|
  \frac1m\log M_{\widetilde A}^{(r)}[k]
  -\frac1m\log M_A^{(r)}[k]\right|
  \le \frac{\Delta_m^{\rm diag}}{m},
  \qquad
  \left|
  \frac1m\log \Lambda_{\widetilde A}^{(r)}[k]
  -\frac1m\log \Lambda_A^{(r)}[k]\right|
  \le \frac{2\Delta_m^{\rm diag}}{m}.
\]
The first comparison is uniform over \(|T|=m+r\), and the normalized
column-Gram comparison is its square.  Hence these lower sides transfer
whenever \(\Delta_m^{\rm diag}=o_\PP(m)\), and the corresponding lower
envelopes transfer on a common coupling whenever
\(\Delta_m^{\rm diag}=o(m)\) almost surely.  This remains only a diagonal
lower-side transfer; no upper side, exact law, row-Gram statement for \(r>0\),
or non-diagonal column-covariance statement is implied.
\end{remark}

\begin{corollary}[Normalized Gaussian fixed-excess fixed-\(k\) lower sides]\label{cor:normalized-gaussian-fixed-excess-fixed-k}
Let \(\F\in\{\R,\C\}\), and let \(G_m\in\F^{N_m\times m}\) have iid standard
Gaussian entries with \(N_m/m\to\gamma>1\).  Let \(U_m\) and \(V_m\) be the row-
and column-normalized matrices of Corollaries~\ref{cor:row-normalized-gaussian}
and~\ref{cor:column-normalized-gaussian}.  Fix \(r\ge0\) and \(k\ge1\).  For
\(m\ge k\), \(N_m\ge m+r\), and either choice \(B_m=U_m\) or \(B_m=V_m\), put
\[
  M_B^{(r)}[k]
  =
  \min_{\substack{T\subset[N_m]\\ |T|=m+r}}\sigma_{m-k+1}((B_m)_T),
\]
and
\[
  \Lambda_B^{(r)}[k]
  =
  \min_{\substack{T\subset[N_m]\\ |T|=m+r}}
  \lambda_{(k)}\!\left((m+r)^{-1}(B_m)_T^*(B_m)_T\right).
\]
Then, for every \(\varepsilon>0\),
\[
  \PP\left\{M_B^{(r)}[k]\le
  \exp\!\left[-\left(\frac{h(\gamma)}{d_\F k(k+r)}
    +\varepsilon\right)m\right]\right\}
  \to0,
\]
and
\[
  \PP\left\{\Lambda_B^{(r)}[k]\le
  \exp\!\left[-\left(\frac{2h(\gamma)}{d_\F k(k+r)}
    +\varepsilon\right)m\right]\right\}
  \to0.
\]
On any common coupling over \(m\), the corresponding almost-sure lower
envelopes hold:
\[
  \liminf_{m\to\infty}\frac1m\log M_B^{(r)}[k]
  \ge -\frac{h(\gamma)}{d_\F k(k+r)},
  \qquad
  \liminf_{m\to\infty}\frac1m\log \Lambda_B^{(r)}[k]
  \ge -\frac{2h(\gamma)}{d_\F k(k+r)} .
\]
\end{corollary}

\begin{proof}
Gaussian \((m+r)\)-row submatrices have rank \(m\) almost surely.  For
\(B_m=U_m\), the normalization is a positive random row diagonal weighting of
\(G_m\); for \(B_m=V_m\), it is a positive random column diagonal weighting.
The Gaussian norm estimates in the proofs of
Corollaries~\ref{cor:row-normalized-gaussian} and
\ref{cor:column-normalized-gaussian} give the corresponding diagonal logarithmic
extrema \(o_\PP(m)\).  The upper row-norm, upper column-norm, and lower
column-norm tail bounds there are summable; the displayed row ball-volume lower
tail is summable after the union over \(N_m=O(m)\) rows as well.  Hence the same
diagonal logarithmic extrema are \(o(m)\) almost surely on any common coupling.
Apply Remark~\ref{rem:random-diagonal-fixed-k} with \(A_m=G_m\), using
Corollaries~\ref{rect-kth:fixed-aspect-lower} and~\ref{rect-kth:gram-lower} for
the lower sides and Corollary~\ref{cor:as-geometric-lower-envelopes} for the
Gaussian almost-sure lower envelopes.
\end{proof}

\begin{corollary}[Quaternion diagonal and normalized fixed-excess fixed-\(k\) lower sides]\label{cor:quat-diagonal-normalized-fixed-excess-fixed-k}
Let \(G_m\in\HH^{N_m\times m}\) have iid standard quaternion Gaussian entries
with \(N_m/m\to\gamma>1\).  Fix \(r\ge0\) and \(k\ge1\).  Let
\[
  \widetilde G_m=D_m^{\rm row}G_mE_m,
\]
where \(D_m^{\rm row}\) and \(E_m\) are positive real diagonal matrices, random
and possibly dependent on \(G_m\), with strictly positive finite entries, and
where either diagonal factor may be absent.  Assume
\[
  \max_i|\log(D_m^{\rm row})_{ii}|
  +
  \max_j|\log(E_m)_{jj}|
  =o_\PP(m),
\]
again omitting absent terms.  Put
\[
  M_{\widetilde G,\HH}^{(r)}[k]
  =
  \min_{\substack{T\subset[N_m]\\ |T|=m+r}}
  \sigma_{m-k+1}((\widetilde G_m)_T)
\]
and
\[
  \Lambda_{\widetilde G,\HH}^{(r)}[k]
  =
  \min_{\substack{T\subset[N_m]\\ |T|=m+r}}
  \lambda_{(k)}\!\left((m+r)^{-1}(\widetilde G_m)_T^*
  (\widetilde G_m)_T\right),
\]
with quaternionic singular values and Hermitian eigenvalues counted once.  Then,
for every \(\varepsilon>0\),
\[
  \PP\left\{M_{\widetilde G,\HH}^{(r)}[k]\le
  \exp\!\left[-\left(\frac{h(\gamma)}{4k(k+r)}
    +\varepsilon\right)m\right]\right\}
  \to0
\]
and
\[
  \PP\left\{\Lambda_{\widetilde G,\HH}^{(r)}[k]\le
  \exp\!\left[-\left(\frac{h(\gamma)}{2k(k+r)}
    +\varepsilon\right)m\right]\right\}
  \to0.
\]
If the displayed logarithmic extreme is \(o(m)\) almost surely on a common
coupling, then the corresponding a.s. lower envelopes also hold with the same
bases.

In particular, the same conclusions hold when \(D_m^{\rm row}\) row-normalizes
the quaternion Gaussian rows, or when \(E_m\) column-normalizes the full
quaternion Gaussian columns.
For the square least-singular-value case \(r=0,k=1\), the stronger convergence
in probability under the same diagonal hypotheses is
Corollary~\ref{cor:quat-square-diagonal-normalized-exact}.
\end{corollary}

\begin{proof}
The entries of the diagonal factors are positive real scalars, hence central in
\(\HH\), and the min--max singular-value sandwich in
Remark~\ref{rem:random-diagonal-fixed-k} applies with the quaternionic
conventions fixed above.  The source lower sides are
Theorem~\ref{thm:quat-fixed-excess-fixed-k-lower} and
Corollary~\ref{cor:quat-rect-kth-gram-lower}; the same remark gives the
corresponding common-coupling lower envelopes when the logarithmic extreme is
\(o(m)\) almost surely.

For row normalization, the squared norm of one row is \((1/4)\chi^2_{4m}\); for
column normalization, the squared norm of one column is \((1/4)\chi^2_{4N_m}\).
The Markov upper-tail estimates and ball-volume lower-tail estimates used in
Corollaries~\ref{cor:row-normalized-gaussian} and
\ref{cor:column-normalized-gaussian}, with real dimensions \(4m\) and \(4N_m\),
show that the corresponding row or column logarithmic extremes are \(o_\PP(m)\)
and \(o(m)\) almost surely after the \(O(m)\) unions.  Gaussian rows and columns
are nonzero almost surely, so the normalizations are well defined.
\end{proof}

\begin{corollary}[Two-sided diagonal transfer]\label{cor:two-sided-diagonal-transfer}
Let \(G_k\in\F^{N_k\times n_k}\) have iid standard Gaussian entries with
\(N_k/k\to\gamma_R>1\) and \(n_k/k\to\gamma_C>1\).  Let
\[
  A_k=D_kG_kE_k
\]
where \(D_k\) and \(E_k\) are deterministic positive diagonal matrices.  Assume
\[
  \frac1k\log\!\left((\min_i D_{k,ii})(\min_j E_{k,jj})\right)\to0,\qquad
  \frac1k\log\!\left((\max_i D_{k,ii})(\max_j E_{k,jj})\right)\to0.
\]
For
\[
  M_A^{\rm two}=\min_{|R|=|C|=k}\smin((A_k)_{R,C}),\qquad
  M_G^{\rm two}=\min_{|R|=|C|=k}\smin((G_k)_{R,C}),
\]
one has
\[
  \frac1k\log M_A^{\rm two}-\frac1k\log M_G^{\rm two}\to0
\]
deterministically.  Thus the two-sided lower side of
Theorem~\ref{two:lower} transfers to \(A_k\).
\end{corollary}

\begin{proof}
For each \(R,C\), \((A_k)_{R,C}=D_R(G_k)_{R,C}E_C\).  Therefore
\[
  (\min_i D_{k,ii})(\min_j E_{k,jj})\,\smin((G_k)_{R,C})
  \le \smin((A_k)_{R,C})
  \le
  (\max_i D_{k,ii})(\max_j E_{k,jj})\,\smin((G_k)_{R,C}).
\]
Take minima over \(R,C\), logs, and divide by \(k\).
\end{proof}

The same pathwise comparison remains true if the positive diagonal matrices
\(D_k,E_k\) are random and possibly dependent on \(G_k\): the absolute
difference of the normalized logarithms is bounded by
\[
  \frac1k\left(\max_i|\log D_{k,ii}|+\max_j|\log E_{k,jj}|\right).
\]
Consequently the two-sided lower-side transfer holds in probability whenever
this bound is \(o_\PP(1)\), and it holds pathwise for almost-sure lower
envelopes whenever the bound is \(o(1)\) almost surely.  This is still only a
diagonal statement.

\begin{remark}
The two-sided diagonal transfer should not be replaced by a general column
covariance statement.  Column selection does not commute with multiplication by
a non-diagonal \(\Sigma^{1/2}\), since
\((G\Sigma^{1/2})_{R,C}=G_R(\Sigma^{1/2})_{:,C}\), not a selected Gaussian minor
times a \(C\times C\) covariance factor.
\end{remark}

\subsection{Gram and sample-covariance transfer}

At the exponential scales of this paper, passing from singular values to Gram
or sample-covariance eigenvalues is deterministic bookkeeping.  It changes
logarithmic constants by a factor of two and preserves the proof status of the
underlying singular-value statement.

\begin{corollary}[Gram and sample-covariance transfer]\label{cor:gram-transfer}
Let \(B_{m,I}\), \(I\in\mathcal I_m\), be a finite family of \(q_m\times n_m\)
matrices over \(\R\), \(\C\), or \(\HH\), with \(q_m\ge1\).  In the quaternionic
case use the adjoint and singular-value conventions fixed above.  Put
\[
  M_m=\min_{I\in\mathcal I_m}\smin(B_{m,I}),\qquad
  \Lambda_m=\min_{I\in\mathcal I_m}
  \lambda_{\min}\!\left(q_m^{-1}B_{m,I}^*B_{m,I}\right).
\]
Then pathwise
\[
  \Lambda_m=q_m^{-1}M_m^2.
\]
Consequently, if \(a_m\to\infty\), \(\log q_m=o(a_m)\), and
\[
  \frac1{a_m}\log M_m\to L
\]
in probability, then
\[
  \frac1{a_m}\log\Lambda_m\to 2L
\]
in probability.  The same implication holds for deterministic comparisons and
for almost-sure lower envelopes, with a lower bound \(L\) replaced by the lower
bound \(2L\).
\end{corollary}

\begin{proof}
For every index \(I\),
\[
  \lambda_{\min}(B_{m,I}^*B_{m,I})=\smin(B_{m,I})^2.
\]
Since \(x\mapsto q_m^{-1}x^2\) is increasing on \([0,\infty)\), taking the
minimum over \(I\) gives the pathwise identity.  Taking logarithms gives
\[
  \frac1{a_m}\log\Lambda_m
  =
  2\frac1{a_m}\log M_m-\frac{\log q_m}{a_m},
\]
on the event \(M_m>0\), and the stated consequences follow from
\(\log q_m=o(a_m)\).
\end{proof}

For iid standard Gaussian square row submatrices over \(\F\in\{\R,\C\}\),
Corollary~\ref{cor:gram-transfer} and the source theorem give
\[
  \frac1m\log
  \min_{\substack{T\subset[N_m]\\ |T|=m}}
  \lambda_{\min}\!\left(m^{-1}A_{m,T}^*A_{m,T}\right)
  \xrightarrow{\PP}
  -\frac{2h(\gamma)}{d_\F}.
\]
The quaternionic theorem above gives, for iid standard quaternion Gaussian
entries,
\[
  \frac1m\log
  \min_{\substack{T\subset[N_m]\\ |T|=m}}
  \lambda_{\min}\!\left(m^{-1}A_{m,T}^*A_{m,T}\right)
  \xrightarrow{\PP}
  -\frac{h(\gamma)}{2}.
\]
Combining the same identity with Proposition~\ref{prop:cov-transfer}, the
square sample-covariance exponent is unchanged for
\(A_m=D_mG_m\Sigma_m^{1/2}\) under the subexponential conditioning hypotheses
there, in particular for uniformly elliptic Gaussian samples with common
population covariance.  A subexponential condition number alone is not enough:
an exponential scalar multiple of the covariance shifts the logarithmic edge.
Fixed-excess, two-sided, and finite-budget lower sides also translate by the
same identity, but only with their original status.  A lower-side singular-value
theorem gives a lower-side Gram theorem; a conditional exact theorem remains
conditional.

\begin{remark}
This corollary is only a change of variables from singular values to Gram
eigenvalues.  It does not identify the proportional-excess Wishart lower edge
on the \(\sqrt m\) singular-value scale, nor the corresponding \(m\)-scale
sample-covariance edge.  Those require the left-tail large-deviation input
described in Section~\ref{sec:geometric}.
\end{remark}

\subsection{General rows}

For independent rows, entry independence is not needed for the square lower
side.  Corollary~\ref{cor:directional-row-lower} gives a finite budget from
uniform affine directional anti-concentration of the row laws.  The fixed-excess
analogue is Corollary~\ref{rect:directional-row-fixed-excess}: the hard-edge
exponent \(r+1\) comes from \(r+1\) remaining rows tested against the normal to
an \(m-1\)-row block.  Thus iid real isotropic log-concave rows have the square
lower-side base \(h(\gamma)\), by Corollary~\ref{cor:logconcave-lower}, and for
each fixed \(r\ge0\),
\[
  \PP\left\{M_m^{(r)}
  \le \exp\left[-\left(\frac{h(\gamma)}{r+1}+\varepsilon\right)m\right]\right\}
  \to0 .
\]
The exact laws remain open: the second moment would need simultaneous
incompressibility of random normals and local anti-concentration for two
projections of a fresh row, now without product coordinates.

\begin{problem}[Log-concave row-submatrix hard edge]
Let \(X_1,\dots,X_N\) be iid isotropic log-concave vectors in \(\R^m\), with
\(N/m\to\gamma>1\), and let \(A\) be the matrix with rows \(X_i\).  Determine
whether
\[
  \frac{1}{m}\log\min_{|T|=m}\smin(A_T)\to -h(\gamma)
\]
in probability.  The lower bound at every threshold below \(-h(\gamma)\) is
proved above; the open part is the matching upper side and the
pair-correlation mechanism for dependent coordinates.
\end{problem}

\begin{problem}[Non-square sample-covariance edge]\label{prob:sample-cov-edge}
For iid samples \(X_i\in\R^m\), determine the extremal lower edge of
\(\ell^{-1}A_T^*A_T\) over row sets \(|T|=\ell\ge m\), beyond the square
Gaussian and elliptic Gaussian cases covered by
Corollary~\ref{cor:gram-transfer}.  For fixed \(\ell-m=r\), this is the
squared fixed-excess hard-edge problem and has only the transferred lower side
unless Conjecture~\ref{rect:exact} is proved.  For \(\ell/m\to\zeta>1\), the
exponential scale is trivial and the remaining question is the constant-order
Wishart edge governed by Problem~\ref{prob:rectangular}.
\end{problem}
